% paper/sections/07_0_scheme_per_variable.tex
%
% §7.0  変数別スキーム方針（SP-L §3）
%        §7 冒頭に配置し, 後続の §7.1 (CLS 移流) / §7.3 (FCCD 移流) /
%        §7.6 (粘性 3 層) の選択根拠を先取りする．
%
% ═══════════════════════════════════════════════════════════════════════════════
\subsection{変数別スキーム選択方針}
\label{sec:scheme_per_variable}
% ═══════════════════════════════════════════════════════════════════════════════

SP-L \cite{SPL} §3 の中心的主張は「差分演算子は\textbf{場の特性}に合わせて
選択せよ」である．
CCD の $\mathcal{O}(h^6)$ 精度はステンシル支持内で場が $C^\infty$ 滑らかである
仮定に依存するため，$C^0$ kink・$\delta$ 関数・step profile を持つ変数に
一律 CCD を適用すると\textbf{演算子の前提違反}が発散・膨張・スプリアス振動として
表出する．

\noindent\textbf{変数別割り当て}：
\begin{center}
\begin{tabular}{@{}lll@{}}
\hline
変数 & 特性 & 推奨スキーム \\
\hline
$\psi$（CLS）& $C^0$ 急 tanh & WENO5 + DCCD hybrid（\S\ref{sec:cls_advection}） \\
$u, v$ bulk & $C^\infty$ 滑 & CCD / UCCD6 / FCCD（\S\ref{sec:fccd_advection}） \\
$u, v$ 界面帯 & $\nabla u$ kink & WENO5 or 2 次片側（\S\ref{sec:viscous_3layer}）\\
$\phi$（再初期化）& Eikonal & WENO--HJ~\cite{JiangPeng2000}（\S\ref{sec:eikonal_reinit}） \\
$p$ & 界面で jump & Face-flux + GFM/IIM~\cite{Fedkiw1999}（\S\ref{sec:balanced_force}）\\
$\rho, \mu$ & step profile & 低次（\S\ref{sec:viscous_3layer}）\\
\hline
\end{tabular}
\end{center}

\noindent\textbf{設計原則（SP-L §3）}：
\begin{itemize}
  \item \textbf{$\psi$ に CCD を適用してはならない．}
        tanh プロファイルを CCD で微分すると急勾配帯の曲率が増幅され，
        CN 半陰的処理でも不安定化する（CHK-137 検証）．
        DCCD スペクトルフィルタ（\S\ref{sec:dccd_derivation}）で散逸を加えた
        場合のみ実用的．
  \item \textbf{$p$ の CCD 節点値計算は jump 補正なしでは禁止．}
        圧力跳躍を滑らかに内挿するため節点 CCD は実効的に
        全域 $\mathcal{O}(1)$ のスプリアス圧力勾配を生む．
        面フラックス + GFM/IIM，あるいは FCCD の面中心演算子
        （\S\ref{sec:fccd_def}）を用いる．
  \item \textbf{$\rho, \mu$ に CCD を適用してはならない．}
        step プロファイルを CCD で差分すると Gibbs 振動が生じ，
        物性値が $[\rho_\text{min},\rho_\text{max}]$ 範囲外になる．
        物性値は常に低次（算術平均 / 調和平均 / $\psi$ 線形内挿）で評価．
  \item \textbf{bulk 速度は CCD 推奨．}
        $u, v$ は $C^\infty$ 滑らか（界面帯を除く）．
        純 CCD でも安定だが，移流のスペクトル散逸欠如を補う
        DCCD または UCCD6（\S\ref{sec:uccd6_def}）が製品版推奨．
  \item \textbf{界面帯の速度は低次でよい．}
        $\nabla u$ の kink を CCD stencil が跨ぐと $\mathcal{O}(1)$ 振動．
        Definition $|\phi|\le 3h$ or phase-indicator mismatch で帯を同定し，
        2 次中心 / 2 次片側に落とす（SP-K §6）．
\end{itemize}

\noindent\textbf{速度-PPE 整合性}（SP-L §2）：
CLS 移流（Stage A, \S\ref{sec:cls_stages}）と運動量移流
（\S\ref{sec:fccd_advection}）は\textbf{直前の PPE 投影済み $\mathbf{u}^n$}
を用いる．
未投影 $\mathbf{u}^*$ を使うと $\mathbf{u}^*$ の発散残差が $\psi$ 質量変化と
PPE RHS の不整合を同時に引き起こす．
正しい順序は：
\begin{equation*}
  \text{PPE}_{n}\!\to\!\mathbf{u}^n\!\to\!\text{CLS 移流} \& \text{運動量移流}
  \!\to\!\mathbf{u}^*_{n+1}\!\to\!\text{PPE}_{n+1}
\end{equation*}
（詳細は \S\ref{sec:algorithm}）．

\noindent\textbf{本章の構成}：
\S\ref{sec:advection_motivation} で CLS 移流 DCCD（既存），
\S\ref{sec:fccd_advection} で FCCD 移流 Option B/C（SP-D 新規），
\S\ref{sec:cls_compression} で CLS 再初期化（既存），
\S\ref{sec:cls_stages} で CLS A--F 6 段階（SP-L §5--§6 新規），
\S\ref{sec:viscous_3layer} で粘性 3 層アーキテクチャ（SP-K §3--§5 新規）．

% ═══════════════════════════════════════════════════════════════════════════════
